\section{Project Contribution}

This sections desribes who contributed what to the project.

\subsection*{Moritz Hörmansdorfer}
Moritz Hörmansdorfer was heavily involved in the conception and development of the modeling and data processing pipelines. His core responsibilities included exploring modeling approaches, defining the data pipeline, as well as training and evaluating the initial machine learning models. Another main focus of his work was critically examining the data foundation, for example, by investigating the tree cadastre as a "ground truth" and considering whether a single model is sufficient for different climate zones (like Vienna and Barcelona). Furthermore, he was responsible for defining and implementing the final database structure. Like all team members, he actively participated in general project planning, report writing, and preparing presentations.

\subsection*{Nina Salnikow}
Nina Salnikow acted as the primary data manager within the project. She led the exploration, collection, and preparation of various datasets, including tree cadastres, Copernicus data, OpenStreetMap (OSM), and other potential sources for different cities (e.g., Graz, Innsbruck). She handled complex scraping tasks, evaluated the feasibility of utilizing Sentinel-2 data, and explored ways to differentiate between private and public areas. Compiling the testing data also fell under her purview. In addition, she was deeply involved in creating the BEAM (Business Event Analysis and Modelling) project plan and took charge of documenting the official sparring group minutes.

\subsection*{Dominik Oberhumer}
Dominik Oberhumer played a central role in the technical project administration and the structuring of the reports. He set up the Git repository and provided the LaTeX templates for all necessary documentation. In terms of project planning, he focused on communication, role distribution, and drafting the schedule, which included the integration of GANTT charts. Furthermore, he was primarily responsible for the formal structuring of the intermediate and final reports and led the comparison of the team's model with existing literature (e.g., Freudenberg et al.). The organizational tasks regarding the sparring groups and inviting the data coaches to the final presentation were also managed by him.

\subsection*{Sara Klasová}
Sara Klasová focused on literature research, Exploratory Data Analysis (EDA), and systematic model validation. She was responsible for drafting the project plan sections concerning the schedule and relevant literature. A major focus of her work was conducting and continuously adapting the EDA, such as integrating new data from Graz and Innsbruck or comparing tree populations based on city size. She also dealt extensively with methodological challenges, like handling different OSM labels across various cities. In the later phases of the project, her core competencies were applied to improving overall model performance, performing validation, and conducting detailed error analysis.

